\section{The Woodbury Portfolio Solver}\label{sec:algorithm}

\subsection{Active-set reduction and the inner solve}

By the change of variables $y = D^{1/2} w$ of Section~\ref{sec:rmt} the
minimum-variance problem is solved against the cleaned correlation $C_0$: minimise
$y^\top C_0\, y - \rho\,\tilde\mu^\top y$ subject to $\tilde B y = c$, $y \geq 0$,
with $\tilde B = B D^{-1/2}$ and $\tilde\mu = D^{-1/2}\hat\mu$, and the portfolio
recovered as $w = D^{-1/2} y$. To lighten notation the rest of this section writes
$w$, $B$, $\hat\mu$ for the standardised weight $y$, the transformed balance matrix
$\tilde B$, and $\tilde\mu$; every displayed formula is in these standardised
coordinates and the reported weights are $D^{-1/2}$ times the solve's output.

This is the equality-augmented non-negative quadratic
of~\cite{schmelzer2026nncg} at $A = 2C_0$, $b = \rho\hat\mu$. We adopt that
solver as printed and contribute the inner linear algebra
(Section~\ref{sec:woodbury}); this subsection records only what the
instantiation needs.

\paragraph{The restricted subproblem}
Fix a set $\mathcal{A} \subseteq \{1,\ldots,n\}$ of active assets, write
$B_\mathcal{A}$ for the columns of $B$ indexed by $\mathcal{A}$, and put
$n_a = |\mathcal{A}|$. The balance-constrained subproblem on $\mathcal{A}$,
\[
  \min_{w_\mathcal{A}:\,B_\mathcal{A} w_\mathcal{A}=c}\;
  w_\mathcal{A}^\top C_{0,\mathcal{A}}\, w_\mathcal{A}
  - \rho\,\hat\mu_\mathcal{A}^\top w_\mathcal{A}\,,
\]
is a saddle-point system whose $p$ multipliers are eliminated analytically by a
$p \times p$ Schur complement~\cite[Section~3]{schmelzer2026nncg}, leaving
\[
  w^\star_\mathcal{A}
  \;=\; \tfrac12 V_1\,\lambda + \tfrac{\rho}{2} v_2\,,
  \qquad
  V_1 = (C_{0,\mathcal{A}})^{-1}B_\mathcal{A}^\top\,,
  \quad
  v_2 = (C_{0,\mathcal{A}})^{-1}\hat\mu_\mathcal{A}\,,
\]
with $\lambda \in \mathbb{R}^p$ fixed by feasibility through
\[
  S\,\lambda = 2c - \rho\,B_\mathcal{A} v_2\,,
  \qquad
  S = B_\mathcal{A} V_1 = B_\mathcal{A}(C_{0,\mathcal{A}})^{-1}B_\mathcal{A}^\top \succ 0\,.
\]
The factors of two carry the convention $w^\top C_0 w$ in place of
$\tfrac12 w^\top A w$; $S$ is SPD because $C_{0,\mathcal{A}} \succ 0$ and
$B_\mathcal{A}$ has full row rank, and the dense $p \times p$ solve is negligible
for the small $p$ typical of balance conditions. In the budget case ($p = 1$,
$B_\mathcal{A} = \mathbf{1}^\top$, $c = 1$) this collapses to a single vector
$v_1 = (C_{0,\mathcal{A}})^{-1}\mathbf{1}$, a scalar Schur complement
$S = \mathbf{1}^\top v_1$, and
$w^\star_\mathcal{A} = \theta v_1 + \tfrac{\rho}{2} v_2$ with
$\theta = \tfrac{\lambda}{2} = (1 - \tfrac{\rho}{2}\mathbf{1}^\top v_2)/(\mathbf{1}^\top v_1)$.
Every outer step thus reduces to $p+1$ applications of
$(C_{0,\mathcal{A}})^{-1}$.

\paragraph{The outer loop}
Non-negativity is enforced by~\cite[Algorithm~1]{schmelzer2026nncg} unchanged,
and its Theorem~4.1 gives finite termination at the unique global minimiser with
no non-degeneracy assumption. The hypotheses hold for \emph{every} active set
here, since $C_{0,\mathcal{A}} \succ 0$ whenever $\bar\mu > 0$ and all
$\delta_j > 0$. That theorem also licenses any \emph{direct} factorisation of the
free-set system in place of its conjugate-gradient inner solve, so the Woodbury
solve of Section~\ref{sec:woodbury} inherits termination and the KKT certificate
verbatim, with no inexactness condition to verify.

The one problem-specific ingredient the loop needs is the gradient
$g = 2\,C_0 w - \rho\,\hat\mu$ supplying its dual test, evaluated through the
low-rank form $C_0w = \bar\mu w + U_k(\Delta_k(U_k^\top w))$ at cost $O(nk)$
(Proposition~\ref{prop:rmt_spectrum}). The return matrix $X$ never enters the
solve, which runs entirely on the $O(nk)$ factors $(U_k, \Delta_k, \bar\mu)$.

\begin{remark}[Implementation safeguards]\label{rem:safeguards}
The reference implementation runs the batch exchange as an unguarded fast path
(assets with $w_i < -10\,\varepsilon$ are dropped together), with an iteration cap
and an active-set-change check in place of the patience-guarded least-index
fallback that~\cite[Algorithm~1]{schmelzer2026nncg} makes the proof rest on.
Neither safeguard triggered in any experiment of Section~\ref{sec:results}, and
the converged weights match the KKT-Cholesky reference to machine precision
(Section~\ref{ssec:solvers}); a deployment wanting the unconditional guarantee
should use the fallback as printed there.
\end{remark}

\begin{remark}[Tolerance scale]\label{rem:tolerance}
The tolerance $\varepsilon$ applies to the standardised weights (primal step) and
the reduced gradient (dual step). In the standardised coordinates both live on the
$O(1)$ correlation scale, with noise floor $\bar\mu \approx 0.5$, so the default
$\varepsilon = 10^{-6}$ is a standard KKT tolerance for active-set
methods~\cite{nocedal2006}; the change of variables removes the dependence on the
raw return scale that a covariance formulation would carry.
\end{remark}

\subsection{Woodbury direct solve}\label{sec:woodbury}

The system matrix on the active set $\mathcal{A}$ is
$C_{0,\mathcal{A}} = \bar\mu I_{n_a} + U_{k,\mathcal{A}}\Delta_k U_{k,\mathcal{A}}^\top$,
where $U_{k,\mathcal{A}} = U_k[\mathcal{A},:] \in \mathbb{R}^{n_a \times k}$ is the restriction
of the signal eigenvectors to the active assets. The Woodbury matrix
identity~\cite{golub2013}, applied with
$A = \bar\mu I_{n_a}$, $U = U_{k,\mathcal{A}}$, $C = \Delta_k$,
$V = U_{k,\mathcal{A}}^\top$, gives

\[
  \bigl(C_{0,\mathcal{A}}\bigr)^{-1}\!b
  \;=\;
  \frac{b}{\bar\mu}
  \;-\;
  \frac{1}{\bar\mu^2}\,
  U_{k,\mathcal{A}}\,
  W^{-1}
  U_{k,\mathcal{A}}^\top b\,,
  \qquad
  W = \Delta_k^{-1} + \frac{1}{\bar\mu}\,U_{k,\mathcal{A}}^\top U_{k,\mathcal{A}}
  \;\in \mathbb{R}^{k \times k}\,.
\]

Note that after restricting to $\mathcal{A}$ the columns of $U_{k,\mathcal{A}}$
are generally no longer orthonormal: $U_{k,\mathcal{A}}^\top U_{k,\mathcal{A}} \neq I_k$.
The correction matrix $W$ accounts for this; its Cholesky factorisation is used
to evaluate $W^{-1}U_{k,\mathcal{A}}^\top b$. The $p+1$ right-hand sides of the
restricted subproblem all share the single factorisation of $W$, so each
right-hand side beyond the first costs only $O(n_a k)$.

\begin{proposition}[Woodbury solve complexity]\label{prop:woodbury_cost}
One application of $(C_{0,\mathcal{A}})^{-1}$ costs:
\begin{itemize}
  \item $O(n_a k)$ for the two matrix-vector products $U_{k,\mathcal{A}}^\top b$
        and $U_{k,\mathcal{A}} z$;
  \item $O(n_a k^2)$ to form $W = \Delta_k^{-1} + \bar\mu^{-1} U_{k,\mathcal{A}}^\top U_{k,\mathcal{A}}$
        (a rank-$k$ outer product sum over $n_a$ columns);
  \item $O(k^3)$ for the Cholesky factorisation of $W$.
\end{itemize}
The dominant term is $O(n_a k^2)$; the total per-application cost is $O(n_a k^2 + k^3)$.
\end{proposition}

Algorithm~\ref{alg:woodbury} assembles these pieces into the routine
$f(\mathcal{A})$ that~\cite[Algorithm~1]{schmelzer2026nncg} calls once per outer
step.

\begin{algorithm}[ht]
\caption{Woodbury inner solve $f(\mathcal{A})$ for the RMT-cleaned correlation}\label{alg:woodbury}
\begin{algorithmic}[1]
\Require Signal eigenpairs $(U_k, \Delta_k)$ with $\Delta_k = \Lambda_k - \bar\mu I$;
         bulk mean $\bar\mu$;
         active set $\mathcal{A}$;
         balance system $B \in \mathbb{R}^{p \times n}$, $c \in \mathbb{R}^p$
         \textit{(default $\mathbf{1}^\top,\, 1$)};
         return-tilt $\rho \geq 0$; expected returns $\hat\mu$
\Ensure  Restricted optimum $w_\mathcal{A}$ and multipliers $\lambda \in \mathbb{R}^p$
\State $U_a \leftarrow U_k[\mathcal{A}, :]$ \Comment{$n_a \times k$ restriction}
\State $W \leftarrow \mathrm{diag}(1/\Delta_k) + U_a^\top U_a / \bar\mu$
  \hfill\textit{($O(n_a k^2)$)}
\State $L \leftarrow \mathrm{chol}(W)$;\;
       define $\mathrm{inv}(b) \leftarrow b/\bar\mu - U_a\,(L^\top \backslash (L \backslash U_a^\top b)) / \bar\mu^2$
  \hfill\textit{($O(k^3)$ once)}
\State $V_1 \leftarrow \mathrm{inv}(B_\mathcal{A}^\top)$;\;
       $v_2 \leftarrow \mathrm{inv}(\hat\mu_\mathcal{A})$ if $\rho > 0$
  \hfill\textit{($p{+}1$ solves, $O(n_a k)$ each)}
\State $S \leftarrow B_\mathcal{A} V_1$;\; solve $S\lambda = 2c - \rho\,B_\mathcal{A} v_2$
  \hfill\textit{($p \times p$ Schur, $O(p^3)$)}
\State \Return $w_\mathcal{A} = \tfrac12 V_1\lambda + \tfrac{\rho}{2} v_2$,\; $\lambda$
\end{algorithmic}
\end{algorithm}

The outer loop calls Algorithm~\ref{alg:woodbury}, extends its output by zeros
off $\mathcal{A}$, and forms the reduced gradient $r = g - B^\top\lambda$ at
$O(nk)$ from the returned multipliers; in the budget case
$(B^\top\lambda)_i = \lambda$ reduces to the scalar threshold $\nu = 2\theta$.

\subsection{Warm-starting}

\emph{Warm-starting profits from active-set continuity, not initial guesses}.
A direct solver has no initial guess; warm-starting instead passes the previous
active set and weights, cutting the number of outer steps --- the direct-solver
half of the warm-start analysis
of~\cite[Proposition~4.2]{schmelzer2026nncg}. Along the frontier the
Markowitz solution is piecewise-linear in $\rho$ with generically single-asset
breakpoint changes~\cite{markowitz1956}, so adjacent solves on a fine grid share
all but a few active assets and the warm solver typically terminates in one to two
outer steps; this is a structural property of parametric quadratic programming, not
an artefact of the data.

\begin{remark}[Relation to the critical line algorithm]\label{rem:cla}
The same piecewise-linear structure underlies the critical line algorithm
(CLA)~\cite{markowitz1956,niedermayer2010,cvxcla}. On a fixed active set the
restricted solution and multipliers are affine in $\rho$
($w_\mathcal{A}(\rho) = a + \rho b$, with $a = V_1 S^{-1} c$ and $b$ read off the
same Schur solve), so every primal exit and dual entry is the root of an affine
function and the exact frontier costs one Woodbury solve per breakpoint
(Section~\ref{ssec:cla}). CLA enumerates the exact frontier for a \emph{fixed}
covariance, whereas the warm-started grid sweep extends unchanged to the
practitioner's setting where consecutive solves differ in the data, not only in
$\rho$.
\end{remark}

\begin{remark}[Ridge terms and other structured modifications]\label{rem:ridge}
The Woodbury structure survives several practically relevant modifications.
A ridge penalty $\gamma\norm{w}^2$ simply shifts the identity coefficient,
$\bar\mu \mapsto \bar\mu + \gamma$ --- the scalar-identity special case of the
regularising split of~\cite[Section~5]{schmelzer2026nncg}, here available in
closed form. A diagonal loading
$D + U_k\Delta_kU_k^\top$ with non-scalar $D \succ 0$ replaces $b/\bar\mu$
by $D^{-1}b$ at unchanged cost. Further \emph{linear equality} constraints
beyond the balance system (e.g.\ sector-neutrality) are absorbed by appending
rows to $B$; the $p \times p$ Schur complement grows with them and the solve
cost stays linear in $p$. What is
\emph{not} supported is a full-rank perturbation of the bulk spectrum, as
produced by nonlinear shrinkage (Remark~\ref{rem:rie}).
\end{remark}

\subsection{Complexity and crossover analysis}

\begin{proposition}[Woodbury vs.\ KKT-Cholesky crossover]\label{prop:crossover}
KKT-Cholesky assembles $C_{0,\mathcal{A}}$ by a rank-$k$ update and factors it
($n_a^2 k + n_a^3/3$ flops); Woodbury forms and factors the $k \times k$
correction $W$ ($n_a k^2 + k^3/3$ flops). Writing $x = k/n_a$, the Woodbury solve
is cheaper per outer step iff
$n_a k^2 + \tfrac{k^3}{3} < n_a^2 k + \tfrac{n_a^3}{3}$, which (dividing by
$n_a^3$) factors as $(x-1)(x^2 + 4x + 1) < 0$, i.e.\ iff $x < 1$ ($k < n_a$).
\end{proposition}

\begin{remark}[Crossover in practice]\label{rem:crossover}
The crossover condition $k < n_a$ holds across essentially all of the frontier:
on S\&P~500 data $k = 17$ against $n_a = 46$ at the minimum-variance solution
($x \approx 0.37$), and on the synthetic frontier $x = k/n_a$ runs from about
$0.33$ at the high-return end ($n_a = 64$) to about $1$ at the tightest
(minimum-variance) point, where $n_a$ dips to $\approx k = 21$. Only at that
extreme are the two solves comparable; elsewhere the flop-count advantage of
Proposition~\ref{prop:crossover} applies, and the measured per-step gap is larger
still ($5.7\times$ warm, Table~\ref{tab:solvers}), consistent with the Woodbury
working set ($k \times k$) fitting in lower cache levels than the assembled
$n_a \times n_a$ system, though we have not instrumented the memory hierarchy
directly. When $n_a \leq k$ the solver falls back to the KKT-Cholesky branch; the
Woodbury solve itself remains valid for all $n_a \geq 1$, since
$W = \Delta_k^{-1} + \bar\mu^{-1}U_{k,\mathcal{A}}^\top U_{k,\mathcal{A}}$
is the sum of a positive definite diagonal matrix and a positive semidefinite
matrix, hence always invertible.
\end{remark}

\begin{proposition}[Per-solve pipeline complexity]\label{prop:pipeline}
With $s$ outer active-set steps, one solve costs $s\cdot O(n_a k^2 + k^3 + nk)$
(the $O(nk)$ term is the dual-step gradient). For a sweep of $S$ solves sharing the
eigenpairs, the pipeline cost is $O(nTk) + S\,s\,O(n_a k^2 + k^3 + nk)$, so the
$O(nTk)$ randomised-SVD preprocessing (Proposition~\ref{prop:rsvd}) amortises to
$O(nTk/S)$ per solve.
\end{proposition}

At the $n=500$ benchmark the preprocessing share of a 50-point sweep is
$30$--$45\%$; for a single solve it dominates. In every case the pipeline total
stays well below the KKT-Cholesky solve alone at large $n$
(Section~\ref{ssec:scaling}).
